% =============================================================================================
\chapter{Méthodologie et architecture}\label{chap:methodo}

% ---------------------------------------------------------------------------------------------
\section{Architecture du système}

La figure~\ref{fig:chaine} présente la chaîne de décision. À gauche, les données réelles et le modèle de risque ;
à droite, les paramètres sectoriels et la simulation ; les deux se rejoignent dans une table de huit critères par
projet, qui alimente le filtre, la méthode multicritère, l'arbitrage budgétaire et les agents de rédaction, avant la
décision humaine.

\begin{figure}[htbp]
\centering
\begin{tikzpicture}[
  font=\small,
  box/.style={draw=eccgris,rounded corners=3pt,align=center,minimum height=1.05cm,text width=5.6cm,fill=white},
  key/.style={box,draw=eccvert,line width=1pt,fill=eccclair},
  arr/.style={-{Stealth[length=2.2mm]},eccgris,thick},
  node distance=0.55cm]
  \node[box] (sba) {\textbf{Données SBA 7(a)}\\ 910\,033 prêts réels, 2000-2016};
  \node[box,right=0.9cm of sba] (dam) {\textbf{Paramètres Damodaran}\\ coût du capital, dette, par secteur};
  \node[box,below=of sba] (ml) {\textbf{Modèle de risque}\\ EBM calibré : probabilité de défaut};
  \node[box,below=of dam] (sim) {\textbf{Simulation des 40 projets}\\ VAN, délai, DSCR, levier, garanties, ESG};
  \coordinate (milieu) at ($(ml.south)!0.5!(sim.south)$);
  \node[box,below=0.55cm of milieu,text width=12.2cm] (scores) {\textbf{Table des critères}\\
        8 critères pour chacun des 40 projets (risque issu du ML + 7 indicateurs)};
  \node[box,below=of scores,text width=12.2cm] (filtre) {\textbf{Agent de conformité : filtre éliminatoire}\\
        DSCR $< 1$ ou probabilité de défaut $> 7{,}1\,\%$ : Refuser};
  \node[key,below=of filtre,text width=12.2cm] (mcda) {\textbf{Moteur MCDA (déterministe)}\\
        poids AHP ; ELECTRE III (classement) ; ELECTRE TRI (catégorie) ; PROMETHEE II (vérification)};
  \node[box,below=of mcda,xshift=-3.3cm,text width=5.6cm] (rob) {\textbf{Agent de robustesse}\\
        200 scénarios de paramètres, 20 tirages};
  \node[box,below=of mcda,xshift=3.3cm,text width=5.6cm] (arb) {\textbf{Agent d'arbitrage}\\
        budget, bornes, plafond sectoriel};
  \node[box,below=2.2cm of mcda,text width=12.2cm] (llm) {\textbf{Agents fondés sur un LLM}\\
        agent rédacteur (note par projet) et conseiller (questions, outils) ; aucun calcul};
  \node[key,below=of llm,text width=7cm] (comite) {\textbf{Comité d'investissement}\\ fixe les paramètres, décide};
  \draw[arr] (sba) -- (ml);
  \draw[arr] (dam) -- (sim);
  \draw[arr,dashed] (sba.east) -- ++(0.45,0) |- node[pos=0.25,right,font=\scriptsize,text=eccgris]{40 prêts du test} (sim.west);
  \draw[arr] (ml) -- (ml|-scores.north);
  \draw[arr] (sim) -- (sim|-scores.north);
  \draw[arr] (scores) -- (filtre);
  \draw[arr] (filtre) -- (mcda);
  \draw[arr] (mcda.south-|rob.north) -- (rob);
  \draw[arr] (mcda.south-|arb.north) -- (arb);
  \draw[arr] (rob) -- (rob|-llm.north);
  \draw[arr] (arb) -- (arb|-llm.north);
  \draw[arr] (llm) -- (comite);
\end{tikzpicture}
\caption{Chaîne de décision : le ML estime le risque, le MCDA classe et catégorise, l'arbitrage respecte le budget,
les agents LLM rédigent, le comité décide.}
\label{fig:chaine}
\end{figure}

\subsection{De l'architecture proposée à l'architecture réalisée}

Le sujet proposait une architecture à neuf composantes. Le tableau~\ref{tab:archi} indique, pour chacune, ce qui a
été réalisé et les écarts. Conformément au sujet, le moteur MCDA est un outil déterministe et non un agent : seuls
la rédaction et le dialogue reposent sur un modèle de langage.

\begin{table}[htbp]
\centering\footnotesize
\caption{Architecture proposée par le sujet et réalisation}\label{tab:archi}
\begin{tabularx}{\textwidth}{>{\raggedright\arraybackslash}p{3.4cm}X>{\raggedright\arraybackslash}p{2.1cm}}
\toprule
Composante proposée & Réalisation dans le projet & Statut \\
\midrule
1. Agent extracteur (lecture de dossiers) & Non réalisé : nos dossiers sont des données structurées. L'import d'un
fichier CSV ou Excel et le calcul automatique de la probabilité de défaut jouent ce rôle. & Écart assumé \\
2. Agent de conformité & Filtre éliminatoire à règles (DSCR, risque), configurable. & Réalisé \\
3. Agents d'analyse par dimension & Modèle de risque (EBM) et facteurs de risque par projet ; VAN, délai et DSCR
calculés par la simulation. Le TRI et la détection d'anomalies ne sont pas calculés. & Réalisé en partie \\
4. Validation par les analystes & Page de configuration : critères, poids, jugements AHP, profils et filtre
modifiables avant le calcul. & Réalisé \\
5. Moteur MCDA & AHP, ELECTRE III, ELECTRE TRI, PROMETHEE II, implémentés en Python et testés à la main. & Réalisé \\
6. Agent d'arbitrage & Sélection sous budget, bornes de montant et plafond sectoriel (deux règles). & Réalisé \\
7. Agent de robustesse & 200 scénarios de paramètres, 20 tirages de simulation, scénarios « et si » du conseiller. &
Réalisé \\
8. Agent rédacteur & Note par projet par LLM, avec vérification des nombres et note de repli sans LLM. & Réalisé \\
9. Comité d'investissement & Paramètres fixés en amont dans la configuration ; décision humaine en aval. & Réalisé \\
\bottomrule
\end{tabularx}
\end{table}

% ---------------------------------------------------------------------------------------------
\section{Modèle de risque}

\subsection{Protocole}

\begin{enumerate}
  \item Les trois modèles sont entraînés sur 2000-2012 avec les mêmes 13 variables.
  \item Leurs hyperparamètres sont choisis sur l'AUROC de validation, par une grille volontairement courte
        (tableau~\ref{tab:hyper}) : quatre valeurs de $C$ pour la logistique, deux nombres d'interactions pour l'EBM,
        deux nombres de feuilles pour LightGBM, dont le nombre d'arbres est fixé par arrêt anticipé (100 tours sans
        amélioration).
  \item Chaque modèle est calibré par régression isotonique ajustée sur la validation. Pour mesurer honnêtement la
        calibration sur la validation elle-même, on utilise une calibration croisée en cinq blocs.
  \item Le seuil F1 est choisi sur la validation.
  \item Le modèle retenu est choisi sur la validation, selon une règle fixée \emph{avant} d'entraîner LightGBM.
  \item Le test (2015-2016) est lu une seule fois, après le choix, pour les métriques finales.
\end{enumerate}

\begin{table}[htbp]
\centering
\footnotesize\setlength{\tabcolsep}{4pt}
\caption{Réglage des hyperparamètres sur la validation}\label{tab:hyper}
\begin{tabular}{lllrc}
\toprule
Modèle & Hyperparamètres & AUROC val. (brut) & Durée (s) & Retenu \\
\midrule
Logistique & C=0.01 & 0,7348 & 20 & oui \\
Logistique & C=0.1 & 0,7308 & 26 &  \\
Logistique & C=1 & 0,7297 & 26 &  \\
Logistique & C=10 & 0,7293 & 30 &  \\
EBM & interactions=3, taux d'apprentissage=0.05 & 0,7676 & 191 &  \\
EBM & interactions=10, taux d'apprentissage=0.05 & 0,7692 & 279 & oui \\
LightGBM & feuilles=31, arbres=334 & 0,7933 & 18 &  \\
LightGBM & feuilles=127, arbres=147 & 0,7959 & 17 & oui \\
\bottomrule
\end{tabular}
\end{table}


\subsection{Règle de choix du modèle}

\begin{encadre}[Règle de choix]
Le modèle utilisé est celui dont le score de Brier calibré est le plus bas sur la validation, \textbf{parmi les
modèles interprétables} (régression logistique et EBM). Le Brier est retenu parce que le MCDA utilise la valeur de la
probabilité ; l'interprétabilité est exigée parce qu'une banque doit pouvoir expliquer un refus de crédit.
LightGBM n'est pas candidat : il sert d'étalon pour mesurer ce que coûte l'interprétabilité.
\end{encadre}

% ---------------------------------------------------------------------------------------------
\section{Simulation des 40 projets candidats}\label{sec:simulation}

\subsection{Principe}

Les 40 projets sont de vrais prêts tirés au hasard dans le bloc de test (graine fixe 2026), que le modèle n'a jamais
vus. Leur probabilité de défaut est celle du modèle ; leurs indicateurs financiers sont simulés. La vraie issue du
prêt est retirée du tableau avant le tirage, si bien que la simulation ne peut pas l'utiliser. Le nombre de 40 permet
de voir l'effet du veto ($40 \times 39 = 1\,560$ comparaisons) tout en restant lisible.

\subsection{Étapes}

\begin{enumerate}
  \item \textbf{Secteur} : chaque code NAICS à deux chiffres est rattaché à un secteur de Damodaran par une table
        faite à la main (tableau~\ref{tab:damodaran} en annexe) ; les cas flous prennent la moyenne du marché.
  \item \textbf{Paramètres} : le taux d'actualisation est le coût du capital du secteur, tiré à $\pm$\,20\,\% ; la
        croissance des flux est une hypothèse d'équipe par secteur, elle aussi tirée à $\pm$\,20\,\%.
  \item \textbf{Financement} : la dette représente entre 50 et 80\,\% de l'investissement ; le levier dette sur fonds
        propres est donc compris entre 1 et 4.
  \item \textbf{DSCR piloté} : 8 projets \og faibles \fg{} (DSCR minimum tiré entre 0,8 et 1,0), 24 \og moyens \fg{}
        (1,0 à 1,5) et 8 \og solides \fg{} (1,5 à 2,0), soit 20\,/\,60\,/\,20\,\%. Les flux annuels croissent avec un
        bruit de 5\,\%, puis sont mis à l'échelle pour que le pire DSCR tombe exactement sur la valeur tirée.
  \item \textbf{VAN et délai} : dix années de flux explicites, une valeur résiduelle en perpétuité sans croissance
        réduite de moitié par prudence ; la durée du prêt est la durée standard de sa tranche (5, 10 ou 25 ans).
  \item \textbf{Garanties} : couverture nulle si le prêt réel n'avait pas de garantie, sinon tirée entre 0,5 et 1,5.
  \item \textbf{ESG} : une classe A, B ou C tirée selon le groupe de secteurs (industrie plutôt A ou B, services plutôt
        C), puis une note autour de 30, 60 ou 85 avec un bruit de 5 points.
  \item \textbf{Emplois par million} : nombre réel d'emplois soutenus divisé par le montant en millions de dollars.
\end{enumerate}

\begin{encadre}[Ce qui est sourcé et ce qui est un choix]
Le coût du capital (Damodaran) et les repères de DSCR (1,25 ; 1,3 à 1,4 dans l'énergie~\cite{nrel2024}) sont
sourcés. Les autres paramètres (croissance, horizon, décote de la valeur résiduelle, bruit, bornes de dette et de
garantie, centres ESG, durées standard) sont des hypothèses d'équipe, testées en sensibilité. Chaque projet reçoit
enfin un nom lisible et fictif (« Cabinet dentaire Baker ») : l'activité vient du vrai code NAICS, le nom de famille
d'une liste fixe ; aucun vrai nom d'emprunteur n'est utilisé.
\end{encadre}

% ---------------------------------------------------------------------------------------------
\section{Paramétrage de la méthode multicritère}

\subsection{Critères, poids et profils}

Le tableau~\ref{tab:criteres} donne les huit critères, leur sens, leurs poids et les profils d'ELECTRE TRI. Les poids
cibles traduisent l'ordre de priorité d'une banque : d'abord le remboursement (risque, DSCR), puis la rentabilité,
puis la protection (garanties, levier), enfin l'impact (ESG, emplois). Les 28 jugements de l'AHP
(tableau~\ref{tab:ahp} en annexe) ont été initialisés en arrondissant à l'échelle de Saaty les rapports de ces poids
cibles ; ils sont modifiables dans l'application.

\begin{table}[htbp]
\centering
\footnotesize\setlength{\tabcolsep}{4pt}
\caption{Critères, poids AHP et profils ELECTRE TRI ($b_1$ : Refuser $|$ Conditionner ; $b_2$ : Conditionner $|$ Financer)}\label{tab:criteres}
\begin{tabular}{lcccccl}
\toprule
Critère & Sens & Poids cible & Poids AHP & Profil $b_1$ & Profil $b_2$ & Type \\
\midrule
Probabilité de défaut & min & 25\,\% & 25,8\,\% & 5,4\,\% & 3,6\,\% & absolu \\
VAN & max & 15\,\% & 16,1\,\% & 0\,\$ & 85\,353\,\$ & mixte \\
Délai de récupération & min & 8\,\% & 8,2\,\% & 10,0\,ans & 6,5\,ans & mixte \\
DSCR minimum & max & 20\,\% & 20,1\,\% & 1,00 & 1,25 & absolu \\
Dette / fonds propres & min & 10\,\% & 9,4\,\% & 3,00 & 1,50 & absolu \\
Couverture par les garanties & max & 10\,\% & 9,4\,\% & 0,50 & 1,00 & absolu \\
Score ESG & max & 7\,\% & 6,2\,\% & 50 & 70 & absolu \\
Emplois par million \$ & max & 5\,\% & 5,0\,\% & 15,63 & 54,00 & relatif \\
\bottomrule
\end{tabular}
\end{table}


Les profils sont \emph{absolus} lorsqu'un repère existe (probabilité de défaut rapportée au taux de référence,
DSCR de 1 et 1,25, levier, garanties couvrant la moitié puis la totalité du prêt, ESG), afin qu'une catégorie traduise
un standard de crédit plutôt qu'un rang dans le lot. Ils sont relatifs (tertiles du lot) quand aucun repère externe
n'existe.

\subsection{Filtre éliminatoire}

Avant le classement, un projet est écarté si son DSCR minimum est inférieur à 1 ou si sa probabilité de défaut dépasse
deux fois le taux de défaut de référence. Le veto d'ELECTRE est relatif (il empêche $a$ de surclasser $b$ si $b$ est
bien meilleur sur un critère) et n'élimine pas un projet dans l'absolu : un seuil absolu doit donc être un filtre.

\begin{encadre}[Adaptation du taux de référence]
La règle prévoyait \og 2 fois le taux de défaut moyen de l'entraînement \fg{}, soit $2 \times 13{,}48\,\% = 27\,\%$.
Or la probabilité de défaut est calibrée sur 2013-2014 : aucun des 40 projets ne dépasse 18\,\%, et le filtre ne se
déclenchait jamais. Comparer une probabilité d'une période à un taux d'une autre période étant incohérent, le taux de
référence retenu est celui de la validation (3,57\,\%), soit un seuil de 7,1\,\%.
\end{encadre}

\subsection{Seuils, coupure et robustesse}

Les seuils d'ELECTRE III et de PROMETHEE II sont exprimés en pourcentage de l'étendue observée de chaque critère :
$q = 5\,\%$, $p = 20\,\%$, $v = 50\,\%$. La coupure d'ELECTRE TRI vaut $\lambda = 0{,}6$ et la règle d'affectation est
pessimiste. La robustesse est mesurée de trois façons : un paramètre à la fois (quatre jeux de seuils, six valeurs de
$\lambda$) ; 200 scénarios combinés tirant les poids à $\pm$\,30\,\%, les seuils dans les plages testées et $\lambda$
entre 0,5 et 0,8 (une décision est \og fragile \fg{} si sa catégorie change dans au moins 25\,\% des scénarios) ;
et 20 nouveaux tirages des indicateurs simulés pour les mêmes 40 prêts.

% ---------------------------------------------------------------------------------------------
\section{Agent d'arbitrage}

L'arbitrage reçoit les catégories du MCDA et choisit les projets à financer sous un budget global, des montants
minimum et maximum par projet et une part maximale par secteur. Pour la démonstration : budget de 5\,M\$, montants
de 20\,000\,\$ à 2\,M\$ et 35\,\% du budget au plus par secteur. Les projets \og Financer \fg{} passent d'abord, puis,
si on l'autorise, les projets \og sous conditions \fg{}. Deux règles sont proposées (section~\ref{sec:theorie-budget}) :

\begin{itemize}
  \item \textbf{suivre le classement} (règle par défaut) : on parcourt les projets dans l'ordre de priorité du comité
        (catégorie, puis rang ELECTRE III) et l'on retient chacun s'il respecte encore les trois contraintes. C'est
        l'optimum exact lorsque tout projet mieux classé doit passer avant les autres, et chaque projet écarté reçoit
        un motif précis ;
  \item \textbf{maximiser le nombre} : programme linéaire en nombres entiers qui retient le plus de projets
        \og Financer \fg{}, puis \og sous conditions \fg{}, puis, à nombre égal, les mieux classés.
\end{itemize}

% ---------------------------------------------------------------------------------------------
\section{Agents fondés sur un modèle de langage}

\subsection{Agent rédacteur}

Pour chaque projet, Python rassemble les \og faits \fg{} déjà calculés et formatés : catégorie, rangs, huit critères,
motif d'exclusion, fragilité, facteurs de risque, statut budgétaire. Le LLM reçoit ces faits avec la consigne de les
recopier sans en créer et rédige une note de six à neuf phrases pour le comité. Un vérificateur extrait tous les nombres
de la note : si un seul n'appartient pas aux faits, la note est rejetée et remplacée par une note de gabarit, écrite
sans LLM.

\subsection{Conseiller conversationnel}

Le conseiller répond en français aux questions sur les projets, les résultats, des scénarios et la méthode. Il ne voit
pas les données : il dispose de dix outils, des fonctions Python déterministes décrites par un schéma
(tableau~\ref{tab:outils}), et rédige à partir de leurs résultats. Sous chaque réponse, l'utilisateur voit les outils
appelés et leurs résultats bruts.

\begin{table}[htbp]
\centering\footnotesize
\caption{Outils du conseiller}\label{tab:outils}
\begin{tabularx}{\textwidth}{>{\ttfamily}p{3.6cm}X}
\toprule
\normalfont Outil & Rôle \\
\midrule
capacites & ce que le conseiller sait faire, ses limites et les méthodes réellement utilisées \\
resume\_resultats & vue d'ensemble : catégories, exclusions, poids, meilleurs projets, fragilité \\
lister\_projets & liste filtrée et triée (catégorie, texte, critère) \\
fiche\_projet & tous les faits d'un projet, éventuellement sous un scénario de poids, de coupure ou sans filtre \\
comparer\_projets & deux à six projets critère par critère \\
statistiques & moyenne, médiane, minimum, maximum d'un critère, par groupe \\
simuler\_scenario & relance le MCDA avec d'autres poids, une autre coupure ou sans filtre \\
arbitrage\_budget & sélection sous budget, avec simulation d'un autre budget \\
parametres\_methode & paramètres réels et explication de la méthode \\
donnees\_historiques & données SBA : taux de défaut, audit, métriques et erreurs du modèle \\
\bottomrule
\end{tabularx}
\end{table}

Deux contrôles automatiques complètent l'appel d'outils. Chaque nombre de la réponse doit figurer dans les résultats
des outils consultés, sinon un avertissement s'affiche. Toute méthode citée mais non utilisée (par exemple une variante
d'ELECTRE inventée par le modèle) déclenche aussi un avertissement. Le fournisseur est Google Gemini, dans son offre
gratuite ; l'application bascule automatiquement sur d'autres modèles en cas de saturation, et fonctionne sans LLM
(notes de gabarit, questions rapides) si aucune clé n'est configurée.
